\documentclass{article}
\usepackage[utf8]{inputenc}
\usepackage[T1]{fontenc}
\usepackage[a4paper]{geometry}		% Reduire les marges
\geometry{hmargin=3cm,vmargin=3.5cm}
\usepackage[english]{babel} 
\usepackage{amsfonts}
\usepackage{amsthm}
\usepackage{amsmath}
\usepackage{amssymb}
\usepackage{hyperref}
\usepackage{array}
\usepackage[svgnames]{xcolor}
\usepackage{xpatch}
\usepackage{tikz}
\usepackage{mathrsfs}
%\usepackage{graphicx}
%\usepackage{caption}
\usepackage[svgnames]{xcolor}
\usepackage{eurosym}
\usetikzlibrary{positioning,calc}
\renewcommand{\thesection}{\Roman{section}}
\xpatchcmd{\proof}{\itshape}{\normalfont\proofnameformat}{}{}
\newcommand{\proofnameformat}{\bfseries}
\theoremstyle{definition}
\newtheorem{defi}{Definition}[section]
\newtheorem{theo}{Theorem}[section]
\newtheorem{prop}[theo]{Proposition}
\newtheorem{lem}[theo]{Lemma}
\newtheorem{coro}[theo]{Corollary}
\newtheorem{exe}{Example}[section]
\newtheorem{rem}{Remarque}
\newtheorem{nota}{Notation}[section]
\renewcommand\qedsymbol{$\blacksquare$}
\numberwithin{equation}{section}

\DeclareMathOperator{\R}{\mathbb{R}}
\DeclareMathOperator{\Z}{\mathbb{Z}}
\DeclareMathOperator{\Q}{\mathbb{Q}}
\DeclareMathOperator{\N}{\mathbb{N}}
\DeclareMathOperator{\E}{\mathbb{E}}
\DeclareMathOperator{\C}{\mathcal{C}}

\title{Exo 4, correction}       
\author{Guillaume Woessner}
\date{}                       % sans date : celle du jour de compilation

%\pagestyle{headings}          % Pour mettre des entetes avec les titres
                              % des sections en haut de page
\sloppy                       % Ne pas faire deborder les lignes dans la marge

\begin{document}

\maketitle                    % Faire un titre utilisant les donnees
                              % passees : \title, \author et \date

\begin{abstract}
\begin{center}
Corrigé de quelques exercices
\end{center}
\end{abstract}

%\tableofcontents              % Table des matieres


Dans tout ce document, l'espace de probabilité considéré est $(\Omega,\mathcal{F},(\mathcal{F}_n)_n,\mathbb{P})$.



\paragraph{Exercice 1} 
Soit $(S_n)_{n\in \N}$ une marche aléatoire sur $\Z$, et $\tau:=\tau_1 = \min\{n\in\N~:~S_n=1\}$. Montrez que 
$$\E[S_n\mathbf{1}_{\tau>n}]\nrightarrow 0,$$
en utilisant le théorème d'arrêt, et ensuite en n'utilisant pas le théorème d'arrêt.

\begin{proof}
Si on avait $\E[S_n\mathbf{1}_{\tau>n}]\mapsto 0$, les 2 autres conditions d'application du théorème d'arrêt étant vérifiées d'après le cours, on aurait $1=\E[S_\tau]=\E[S_0]=0$. Ainsi par l'absurde on a montré que $\E[S_n\mathbf{1}_{\tau>n}]\nrightarrow 0$.\\
Sans le théorème d'arrêt, on doit remarquer que $0=\E[S_n]=\E[S_n\mathbf{1}_{\tau>n}]+\E[S_n\mathbf{1}_{\tau\leq n}]$, et faire le calcul suivant
\begin{align*}
\E[S_n\mathbf{1}_{\tau>n}]&= -\E[S_n\mathbf{1}_{\tau\leq n}] = -\sum_{k=0}^n \E[S_n\mathbf{1}_{\tau=k}]\\
&= -\sum_{k=0}^n \E[(S_k + X_{k+1}+\dots+X_n)\mathbf{1}_{\tau=k}] \\
&= -\sum_{k=0}^n \E[S_k \mathbf{1}_{\tau=k}] - \E[(X_{k+1}+\dots+X_n)\mathbf{1}_{\tau=k}]\\
&= -\sum_{k=0}^n \E[1. \mathbf{1}_{\tau=k}] - 0 \qquad \text{par indépendance}\\
&= -\sum_{k=0}^n \mathbb{P}[\tau=k] = -\mathbb{P}(\tau \leq n) \mapsto -1.
\end{align*}
\end{proof}




\paragraph{Exercice 2} 
Montrez que le corollaire suivant du théorème d'arrêt est vrai.\\
Soit $(X_n)_n$ une martingale. Si un temps d'arrêt $\tau$ est presque sûrement fini, et il existe une variable aléatoire intégrable $M$ telle que pour tout $n$,
$$|X_{n\wedge \tau}|\leq M \qquad \textit{as}.$$
Alors $\E[X_\tau]=\E[X_0].$

\begin{rem}
Ce corollaire n'est pas dans le cours, mais il est très proche du point ii) du Corollaire 3.4 du cours.
\end{rem}

\begin{proof}
Il s'agit de montrer les conditions ii) et iii) du théorème d'arrêt.
\begin{itemize}
\item[ii)] On a $\E[|X_T|]=\E[\lim_n |X_{n\wedge \tau}|]= \lim_n \E[|X_{\tau\wedge n}|]\leq \lim_n \E[M]=\E[M]<\infty$, la deuxième égalité étant vraie grâce à l'hypothèse et au théorème de convergence dominée.
\item[iii)] On a $|\E[S_n \mathbf{1}_{\tau\wedge n}]|\leq \E[|S_n| \mathbf{1}_{\tau\wedge n}] = \E[|S_{\tau\wedge n}| \mathbf{1}_{\tau\wedge n}] \leq \E[M\mathbf{1}_{\tau\wedge n}]$. Or $M\mathbf{1}_{\tau\wedge n} \mapsto 0$ par l'Exercice 3, et est dominée par $M$ intégrable par hypothèse. Ainsi par le théorème de convergence dominée à nouveau on a bien que $\E[S_n \mathbf{1}_{\tau\wedge n}] \mapsto 0$.
\end{itemize}
Ainsi le théorème d'arrêt s'applique et le résultat s'ensuit.
\end{proof}



\paragraph{Exercice 2'} 
Montrez que le corollaire suivant du théorème d'arrêt est vrai.\\
Soit $(X_n)_n$ une martingale. Si un temps d'arrêt $\tau$ est intégrable, et il existe un réel $C$ tel que pour tout $n$
$$\E\big[|X_n-X_{n-1}|~\big|~\mathcal{F}_{n-1}\big]\leq C \qquad \textit{as}.$$
Alors $\E[X_\tau]=\E[X_0].$

\begin{rem}
Etant donné que $\E\big[|X_n-X_{n-1}|~\big|~\mathcal{F}_{n-1}\big]\leq C$ est impliqué par $|X_n-X_{n-1}|\leq C$, cet exercice est bien une démonstration du point iii) du Corollaire 3.4 du cours.
\end{rem}

\begin{proof}
On réécrit le processus arrêté $X_{\tau\wedge t}$ par 
$$X_{\tau\wedge n} = X_0 + \sum_{k=1}^{\tau\wedge n} (X_k-X_{k-1}).$$
Ainsi on a $|X_{\tau\wedge n}|\leq M$ pour tout $n$, où 
$$M:=|X_0| + \sum_{k=1}^\tau |X_k-X_{k-1}|=|X_0| + \sum_{k=1}^\infty |X_k-X_{k-1}| \, \mathbf{1}_{\tau+1>k}.$$
Et par le théorème de convergence monotone on obtient 
\begin{align*}
\E[M] &= \E[|X_0|] + \sum_{k=1}^\infty \E[|X_k-X_{k-1}|\,\mathbf{1}_{\tau+1>k}]\\
&= \E[|X_0|] + \sum_{k=1}^\infty \E\big[ \E[ |X_k-X_{k-1}|\,\mathbf{1}_{\tau>k-1}~|~\mathcal{F}_{k-1}] \big]\\
&= \E[|X_0|] + \sum_{k=1}^\infty \E\big[ \E[ |X_k-X_{k-1}|~|~\mathcal{F}_{k-1}] \mathbf{1}_{\tau>k-1} \big]\\
&\leq \E[|X_0|] + \sum_{k=1}^\infty C\E[ \mathbf{1}_{\tau>k+1} ]\\
&= \E[|X_0|] + C\sum_{k=1}^\infty \mathbb{P}[\tau\geq k] = \E[|X_0|] + C\E[\tau] < \infty.
\end{align*}
On obtient ainsi que le processus $|X_{\tau\wedge n}|$ est dominé par une variable $M$ intégrable, et d'après le théorème de convergence dominée on a alors que $\E[X_\tau]=\E[\lim_n X_{\tau\wedge n}] = \lim_n\E[X_{\tau\wedge n}]$.\\
Ainsi, comme on sait que le processus $X_{\tau\wedge n}$ est une martingale, on sait que $\E[X_{\tau\wedge n}]=\E[X_0]$. On peut alors conclure.
\end{proof}





\paragraph{Exercice 4}
Soit $a,b\in\N^\times$, $(S_n)_{n\in \N}$ une marche aléatoire sur $\Z$ et $\tau:=\inf\{n\in \N~:~S_n\notin[-a,b]\}$. Utilisez le fait que $S_n^2-n$ est une martingale (\textit{cf} feuille Exo 2) pour montrer que $\E \tau = ab$. 

\begin{rem}
Je vais donner 2 façons de résoudre cet exercice. La première consiste à montrer la seule chose qui a été admise dans le cours, c'est-à-dire que $\E\tau <\infty$. La seconde est une façon alternative de faire l'exercice, plus simple et indépendante de ce qui a été fait en cours.
\end{rem}

\begin{proof}
On rappelle qu'on a montré en cours que $\tau<\infty$ presque sûrement.\\
\begin{itemize}
\item[\textbf{V1}] Montrons que $\E\tau<\infty$. On définit les évènements $A_n:=\{X_i=1~\forall n(a+b)+1 \leq i \leq (n+1)(a+b)\}$. Les $A_n$ sont des évenements indépendants d'une même probabilité $p>0$, et on définit $\tau' := \inf\{n\geq 0~:~A_n\text{ se réalise}\}$. On a que $\tau\leq (a+b)\tau'$ (voir la preuve du cours que $\mathbb{P}(\tau<\infty)=1$ pour plus de détails). Or $\tau'$ suit une loi géométrique (nombre d'essais avant d'obtenir un succès lors d'une succession d'épreuves de Bernouilli \textit{iid}), et une loi géométrique admet une espérance finie, QED.
\item[\textbf{V2}] Soit $n\in\N^\times$, alors $\tau\wedge n$ est un temps d'arrêt borné, ainsi par le point i) du Corollaire 3.4 du théorème d'arrêt du cours on sait que $\E[S_{\tau\wedge n}^2-\tau\wedge n] = \E[S_0^2-0]= 0$, ainsi $\E[S_{\tau\wedge n}^2]=\E[\tau\wedge n]$.\\
Par le théorème de convergence monotone on sait que $\E[\tau]=\E[\lim_n \tau\wedge n] = \lim_n\E[\tau\wedge n]$.\\
D'autre part, $S_{\tau\wedge n}\mapsto_n S_\tau$ presque sûrement, et $S_{\tau\wedge n}^2 \leq (a\vee b)^2$ pour tout $n$, donc par le théorème de convergence dominée on a également $\E[S_\tau^2]=\E[\lim_n S_{\tau\wedge n}^2]=\lim_n\E[S_{\tau\wedge n}^2]$.\\
Recollons les morceaux, et on obtient $\E[\tau]=\E[S_\tau^2]$.\\
Finalement, il reste à calculer $\E[S_\tau^2]$. Or on a vu en cours que $S_\tau$ est une variable aléatoire prenant la valeur $-a$ avec probabilité $\frac{b}{a+b}$ et la valeur $b$ avec probabilité $\frac{a}{a+b}$. Ainsi un bref calcul permet d'obtenir que $\E[S_\tau^2]=ab$.
\end{itemize}
\end{proof}




\end{document}
